\part{The Algebra of Two-by-Two Matrices}

\chapter{Quaternions and the Group \texorpdfstring{$\SU(2)$}{SU(2)}}
\label{ch:quaternions}

\begin{usesbox}
\textbf{This chapter uses:} linear algebra and the complex numbers, nothing else.\\
\textbf{It provides:} the quaternion algebra $\HH$, its matrix model inside $\Mat_2(\CC)$, the identification of unit quaternions with $\SU(2)$, the rotation action on $\RR^3$, the Hopf map onto the two-sphere, and the decomposition of $\SO(4)$ into left and right multiplications.
\end{usesbox}

Most of physics is written in terms of quadratic forms. The length of a vector, the interval between two events, the energy of a free particle, the Laplacian, the kinetic term of a field: each is a sum of squares, sometimes with signs. A recurring theme of this book is that such forms become far more transparent once they are written as a \emph{product} of two linear factors. The prototype is the factorisation of $x^2+y^2$ as $(x+\ii y)(x-\ii y)$, which is the whole reason complex numbers describe rotations of the plane so efficiently.

Hamilton's discovery of 1843 was that the same trick works in four dimensions, at the price of giving up commutativity \cite{Hamilton1844}. This chapter develops the resulting algebra only as far as we shall need it later. Three facts will be used over and over again: the norm of a quaternion is multiplicative; quaternions are the same thing as a particular real subalgebra of the $2\times 2$ complex matrices; and the unit quaternions are the group $\SU(2)$, which acts on three-space by rotations and on the two-sphere through the Hopf map.

\section{Hamilton's relations}

A quaternion is an expression $q=a+b\,i+c\,j+d\,k$ with real $a,b,c,d$. Addition is componentwise, and multiplication is fixed by bilinearity together with Hamilton's relations
\begin{equation}\label{eq:hamilton}
i^2=j^2=k^2=ijk=-1 .
\end{equation}
From \eqref{eq:hamilton} one derives $ij=k=-ji$, $jk=i=-kj$ and $ki=j=-ik$, so the algebra is associative but not commutative. The real number $a$ is the \emph{real part} $\operatorname{Re}q$ and $b\,i+c\,j+d\,k$ the \emph{imaginary part}; quaternions with vanishing real part are called \emph{pure}. We identify the pure quaternions with $\RR^3$ by $b\,i+c\,j+d\,k\leftrightarrow(b,c,d)$. For two pure quaternions $x$ and $y$ the product packages both of the familiar vector products,
\begin{equation}\label{eq:pureproduct}
xy=-\,x\cdot y+x\times y ,
\end{equation}
which is often the quickest way to remember how quaternions multiply.

\begin{theorem}[Hamilton's relations]\label{thm:hamilton}\ver\
In the quaternion algebra over $\RR$ the elements $i=(0,1,0,0)$, $j=(0,0,1,0)$, $k=(0,0,0,1)$ satisfy $i^2=j^2=k^2=-1$ and $ijk=-1$.
\end{theorem}
\begin{proof}
Write $p=a+b\,i+c\,j+d\,k$ and $q=e+f\,i+g\,j+h\,k$. The product in the quaternion algebra is defined componentwise by
\[
\begin{aligned}
pq={}&(ae-bf-cg-dh)+(af+be+ch-dg)\,i\\
&+(ag-bh+ce+df)\,j+(ah+bg-cf+de)\,k .
\end{aligned}
\]
For $p=q=i$ we have $b=f=1$ and all other components $0$, so only the term $-bf=-1$ survives: $i^2=-1$. In the same way $j^2=-cg=-1$ and $k^2=-dh=-1$. For $p=i$, $q=j$ ($b=1$, $g=1$) the only non-zero term is $bg$ in the $k$-component, so $ij=k$. Then $ijk=(ij)k=k\cdot k=-1$.
\end{proof}


The \emph{conjugate} of $q$ is $\qbar q=a-b\,i-c\,j-d\,k$. Conjugation reverses products, $\qbar{pq}=\qbar q\,\qbar p$, and the key identity is
\begin{equation}\label{eq:qqbar}
q\,\qbar q=\qbar q\,q=a^2+b^2+c^2+d^2=|q|^2 .
\end{equation}
Equation~\eqref{eq:qqbar} is the four-dimensional analogue of $z\bar z=|z|^2$: the Euclidean quadratic form on $\RR^4$ factorises into two linear factors. Every nonzero quaternion therefore has an inverse $q^{-1}=\qbar q/|q|^2$, and $\HH$ is a division algebra. Using \eqref{eq:qqbar} twice gives the multiplicativity of the norm.

\begin{theorem}[Multiplicativity of the norm]\label{thm:normmul}\ver\
For all $p,q\in\HH$, $|pq|^2=|p|^2\,|q|^2$.
\end{theorem}

\begin{proof}
Write $q=a+x$ with $a$ real and $x$ pure. By \eqref{eq:pureproduct}, $x^2=-x\cdot x+x\times x=-|x|^2$, hence
\[
q\qbar q=(a+x)(a-x)=a^2-ax+xa-x^2=a^2+|x|^2=|q|^2,
\]
because the real number $a$ commutes with $x$. This proves \eqref{eq:qqbar}.

Next, conjugation reverses products. It is real-linear, so it suffices to check $\qbar{pq}=\qbar q\,\qbar p$ for $p=a+x$, $q=b+y$. The real parts commute with everything, and for the pure parts \eqref{eq:pureproduct} gives
\[
\qbar{xy}=\overline{-x\cdot y+x\times y}=-x\cdot y-x\times y=-y\cdot x+y\times x=yx=(-y)(-x)=\qbar y\,\qbar x .
\]
Expanding $(a+x)(b+y)$ and conjugating term by term then gives $\qbar{pq}=\qbar q\,\qbar p$.

Finally, using associativity and the fact that the real number $|q|^2$ commutes with $p$,
\[
|pq|^2=(pq)\,\qbar{pq}=p\,q\,\qbar q\,\qbar p=p\,|q|^2\,\qbar p=|q|^2\,p\qbar p=|p|^2|q|^2 .\qedhere
\]
\end{proof}

Written out in components, \cref{thm:normmul} is Euler's four-square identity: a product of two sums of four squares is again a sum of four squares. A consequence that we shall use constantly is that the unit quaternions
\[
S^3=\{q\in\HH:\ |q|=1\}
\]
form a group under multiplication. Topologically this group is the three-sphere.

\section{Quaternions as complex matrices}

Quaternions can be represented faithfully by $2\times2$ complex matrices. For $q=a+b\,i+c\,j+d\,k$ put
\begin{equation}\label{eq:rho}
\rho(q)=\begin{pmatrix} a+b\,\ii & c+d\,\ii\\ -c+d\,\ii & a-b\,\ii\end{pmatrix}.
\end{equation}
Here $\ii$ is the ordinary imaginary unit of $\CC$, which must not be confused with the quaternion $i$. In terms of the Pauli matrices $\sigma_1,\sigma_2,\sigma_3$ one has $\rho(1)=\mathbf 1$, $\rho(i)=\ii\sigma_3$, $\rho(j)=\ii\sigma_2$ and $\rho(k)=\ii\sigma_1$; the three imaginary units are represented by anti-Hermitian, traceless matrices.

\begin{theorem}[The matrix model of $\HH$]\label{thm:rho}\ver\
The map $\rho:\HH\to\Mat_2(\CC)$ of \eqref{eq:rho} has the following properties.
\begin{enumerate}[label=(\roman*)]
\item $\rho(pq)=\rho(p)\rho(q)$ and $\rho(1)=\mathbf 1$.
\item $\rho$ is injective.
\item $\rho(\qbar q)=\rho(q)^\dagger$.
\item $\det\rho(q)=|q|^2$.
\item If $|q|=1$ then $\rho(q)\rho(q)^\dagger=\mathbf 1$ and $\det\rho(q)=1$, that is, $\rho(q)\in\SU(2)$.
\end{enumerate}
\end{theorem}
\begin{proof}
Write $\rho(q)=a\mathbf 1+b\,(\ii\sigma_3)+c\,(\ii\sigma_2)+d\,(\ii\sigma_1)$, which is \eqref{eq:rho} written out.

(i) $\rho$ is real-linear, so it suffices to check products of basis elements. The images $I=\ii\sigma_3$, $J=\ii\sigma_2$, $K=\ii\sigma_1$ satisfy $I^2=-\sigma_3^2=-\mathbf 1$, and likewise $J^2=K^2=-\mathbf 1$. Moreover $\sigma_3\sigma_2=-\ii\sigma_1$, so
\[
IJK=\ii^3\sigma_3\sigma_2\sigma_1=-\ii\,(-\ii\sigma_1)\sigma_1=-\mathbf 1 .
\]
So $I,J,K$ obey Hamilton's relations, the products of basis elements are preserved, and by bilinearity $\rho(pq)=\rho(p)\rho(q)$. Clearly $\rho(1)=\mathbf 1$.

(ii) If $\rho(q)=0$ then $a+b\ii=0$ and $c+d\ii=0$, so $a=b=c=d=0$. Since $\rho$ is linear, $\rho(p)=\rho(q)$ implies $\rho(p-q)=0$, hence $p=q$.

(iii) Taking the conjugate transpose of \eqref{eq:rho},
\[
\rho(q)^\dagger=\begin{pmatrix}a-b\ii&-c-d\ii\\ c-d\ii&a+b\ii\end{pmatrix},
\]
which is \eqref{eq:rho} with $b,c,d$ replaced by $-b,-c,-d$, i.e.\ $\rho(\qbar q)$.

(iv) $\det\rho(q)=(a+b\ii)(a-b\ii)-(c+d\ii)(-c+d\ii)=(a^2+b^2)-(-c^2-d^2)=a^2+b^2+c^2+d^2$.

(v) By (i), (iii) and \eqref{eq:qqbar}, $\rho(q)\rho(q)^\dagger=\rho(q)\rho(\qbar q)=\rho(q\qbar q)=|q|^2\mathbf 1$, which is $\mathbf 1$ for $|q|=1$; and $\det\rho(q)=|q|^2=1$ by (iv).
\end{proof}


Property (iv) is the matrix version of \eqref{eq:qqbar}. Property (v) together with (i) and (ii) identifies the group of unit quaternions with a subgroup of $\SU(2)$; since every element of $\SU(2)$ has the form $\begin{psmallmatrix}\alpha&\beta\\-\bar\beta&\bar\alpha\end{psmallmatrix}$ with $|\alpha|^2+|\beta|^2=1$, which is exactly $\rho$ of a unit quaternion, the subgroup is all of $\SU(2)$:
\[
S^3\;\cong\;\SU(2).
\]
This is the identification that makes quaternions natural for spin. A unit quaternion is a normalised spinor $(\alpha,\beta)\in\CC^2$, written as the first row of a matrix.

\subsection*{Two real forms of the same complex algebra}

The image $\rho(\HH)$ is a real four-dimensional subspace of the complex four-dimensional algebra $\Mat_2(\CC)$. It is the fixed set of the antilinear map $X\mapsto\sigma_2\bar X\sigma_2$. Indeed one checks that $\Mat_2(\CC)=\rho(\HH)\oplus\ii\,\rho(\HH)$, so that $\Mat_2(\CC)$ is the complexification $\HH\otimes_\RR\CC$, the algebra Hamilton called the \emph{biquaternions}. This is the first appearance of a structure that organises the whole of Chapter~2: the same complex algebra $\Mat_2(\CC)$ contains two inequivalent real forms. One of them, studied here, is $\HH$, whose unit sphere is the group of rotations of spin. The other, studied in the next chapter, is the space of Hermitian matrices, which is Minkowski space. Both carry a quadratic form given by the determinant, and in both cases the determinant factorises as a matrix times its adjugate.

\section{Rotations of three-space}

The group $S^3$ acts on $\HH$ by conjugation, $x\mapsto qx\qbar q$. Two short computations show that this action preserves the decomposition into real and pure parts and preserves the norm.

\begin{theorem}[Conjugation preserves the real part and the norm]\label{thm:conj}\ver\
For all $q,x\in\HH$,
\[
\operatorname{Re}(q\,x\,\qbar q)=|q|^2\operatorname{Re}x,\qquad |q\,x\,\qbar q|^2=|q|^4\,|x|^2 .
\]
In particular, for a unit quaternion $q$ the map $x\mapsto qx\qbar q$ sends pure quaternions to pure quaternions and preserves their length.
\end{theorem}
\begin{proof}
For $u=u_0+\mathbf u$ and $v=v_0+\mathbf v$, \eqref{eq:pureproduct} gives $\operatorname{Re}(uv)=u_0v_0-\mathbf u\cdot\mathbf v$, which is symmetric in $u$ and $v$. Hence $\operatorname{Re}(uv)=\operatorname{Re}(vu)$. Applying this with $u=q$ and $v=x\qbar q$,
\[
\operatorname{Re}(q\,x\,\qbar q)=\operatorname{Re}(x\,\qbar q\,q)=\operatorname{Re}(x\,|q|^2)=|q|^2\operatorname{Re}x .
\]
For the norm, \cref{thm:normmul} applied twice gives $|qx\qbar q|^2=|q|^2|x|^2|\qbar q|^2=|q|^4|x|^2$, since $|\qbar q|=|q|$. If $|q|=1$ and $\operatorname{Re}x=0$, then $\operatorname{Re}(qx\qbar q)=0$ and $|qx\qbar q|=|x|$.
\end{proof}


So every unit quaternion gives an orthogonal transformation $R_q$ of $\RR^3$. The map $q\mapsto R_q$ is a homomorphism, and its kernel is $\{\pm1\}$, since only real quaternions commute with all of $i,j,k$. Writing $q=\cos(\theta/2)+\sin(\theta/2)\,\mathbf n$ with $\mathbf n$ a pure unit quaternion, $R_q$ is the rotation by the angle $\theta$ about the axis $\mathbf n$. This is the classical statement that $\SU(2)$ is the double cover of $\SO(3)$; the half angle $\theta/2$ is the reason spinors change sign under a rotation by $2\pi$.

\section{The Hopf map and the Bloch sphere}

Fix the pure unit quaternion $i$ and consider
\begin{equation}\label{eq:hopf}
\pi:S^3\to S^2,\qquad \pi(q)=q\,i\,\qbar q .
\end{equation}
By \cref{thm:conj}, $\pi(q)$ is a pure quaternion of unit length, so $\pi$ does land in the unit sphere of $\RR^3$. Two unit quaternions have the same image exactly when they differ by right multiplication by $e^{i\varphi}=\cos\varphi+i\sin\varphi$, so the fibres of $\pi$ are great circles. This is the Hopf fibration $S^1\to S^3\to S^2$ \cite{Hopf1931}.

\begin{figure}[t]
\centering
\begin{tikzpicture}[scale=1.6]
\draw[thick] (0,0) circle (1);
\draw[dashed] (1,0) arc (0:180:1 and 0.3);
\draw (-1,0) arc (180:360:1 and 0.3);
\draw[-{Stealth}] (0,0) -- (0.55,0.62) node[above right] {$\mathbf n=\psi^\dagger\bs\psi$};
\fill (0,1) circle (0.025) node[above] {$|{\uparrow}\rangle$};
\fill (0,-1) circle (0.025) node[below] {$|{\downarrow}\rangle$};
\draw[dotted] (0.55,0.62) -- (0.55,-0.12);
\node at (1.55,0.2) {$S^2$};
\end{tikzpicture}
\caption{The two-sphere as the space of pure spin states. A normalised spinor $\psi$ (a point of $S^3$) is sent to the unit vector $\mathbf n=\psi^\dagger\bs\psi$; spinors differing by a phase give the same point.}
\label{fig:bloch}
\end{figure}

In the language of quantum mechanics the Hopf map is the passage from a normalised spinor $\psi=(\alpha,\beta)\in\CC^2$ to its Bloch vector $\mathbf n=\psi^\dagger\bs\psi$. The fibre is the overall phase of $\psi$, which is unobservable, so the two-sphere of \cref{fig:bloch} is the space of pure states of a spin-$\frac12$ system, the Bloch sphere \cite{BengtssonZyczkowski}. The same sphere has two further meanings that will matter later. First, the pure unit quaternions are exactly the solutions of $I^2=-1$, and each such $I$ spans a copy $\CC_I=\RR+\RR I$ of the complex numbers inside $\HH$; every quaternion lies in at least one of these complex slices. This is the starting point of the theory of slice-regular functions \cite{GentiliStruppa2007}. Second, in Chapter~2 we shall see that the null directions of Minkowski space form a two-sphere, the celestial sphere, and that it is again the projective line $\CC P^1=S^2$ on which the Lorentz group acts by M\"obius transformations \cite{PenroseRindler1}. That one sphere plays the roles of the imaginary units, the pure spin states and the light rays is not a coincidence of notation. It is the geometric fact on which the first two parts of the book are built.

\section{Left and right multiplication}

Since $\HH$ is not commutative, left and right multiplication are different operators. For $a,b\in\HH$ let $L_a$ and $R_b$ be the real-linear maps $x\mapsto ax$ and $x\mapsto xb$ of $\HH\cong\RR^4$. In the basis $(1,i,j,k)$, for $a=a_0+a_1i+a_2j+a_3k$,
\[
L_a=\begin{pmatrix} a_0&-a_1&-a_2&-a_3\\ a_1&a_0&-a_3&a_2\\ a_2&a_3&a_0&-a_1\\ a_3&-a_2&a_1&a_0\end{pmatrix},\qquad
R_b=\begin{pmatrix} b_0&-b_1&-b_2&-b_3\\ b_1&b_0&b_3&-b_2\\ b_2&-b_3&b_0&b_1\\ b_3&b_2&-b_1&b_0\end{pmatrix}.
\]
By associativity $(ax)b=a(xb)$, so every $L_a$ commutes with every $R_b$. For a pure quaternion $b$ the matrix $R_b$ is antisymmetric, so pure quaternions acting on the right span a copy of $\su(2)\cong\RR^3$ inside $\mathfrak{so}(4)$; the same is true on the left. The next theorem says that these two copies are each other's commutants.

\begin{theorem}[Commutant of the left action]\label{thm:commutant}\ver\
\begin{enumerate}[label=(\roman*)]
\item For a pure quaternion $b$, the matrix $R_b$ is antisymmetric and commutes with $L_i$ and $L_j$.
\item Conversely, every antisymmetric real $4\times4$ matrix $M$ that commutes with $L_i$ and $L_j$ is of the form $R_b$ for a pure quaternion $b$, namely $b=M_{10}\,i+M_{20}\,j+M_{30}\,k$.
\end{enumerate}
\end{theorem}
\begin{proof}
Use the Euclidean inner product $\langle u,v\rangle=\operatorname{Re}(u\,\qbar v)$ on $\HH=\RR^4$.

(i) For all $x,y$, $\langle xb,y\rangle=\operatorname{Re}(xb\,\qbar y)$ and $\langle x,y\qbar b\rangle=\operatorname{Re}(x\,\overline{y\qbar b})=\operatorname{Re}(x\,b\,\qbar y)$. So the transpose of $R_b$ is $R_{\qbar b}$. For pure $b$, $\qbar b=-b$ and therefore $R_b^{\mathsf T}=-R_b$. Associativity gives $L_aR_bx=a(xb)=(ax)b=R_bL_ax$ for every $a$, in particular for $a=i,j$.

(ii) Since $M$ is antisymmetric, $M_{00}=0$, so the first column of $M$ is $Me_0=(0,M_{10},M_{20},M_{30})$, which is the pure quaternion $b$ of the statement. Put $N=M-R_b$. Then $N$ commutes with $L_i$ and $L_j$, hence also with $L_k=L_iL_j$ and with $L_1=\mathbf 1$, and so with every $L_a$ by linearity. Moreover $R_be_0=1\cdot b=b=Me_0$, so $Ne_0=0$. For any quaternion $a$, $a=L_ae_0$, hence
\[
N a=N L_a e_0=L_a N e_0=0 .
\]
So $N=0$, i.e.\ $M=R_b$.
\end{proof}


Because $L_iL_j=L_k$, commuting with $L_i$ and $L_j$ is the same as commuting with all left multiplications. The theorem is the infinitesimal form of the classical isomorphism
\[
\SO(4)\cong\bigl(\SU(2)_L\times\SU(2)_R\bigr)/\{\pm1\},\qquad x\mapsto a\,x\,\qbar b ,
\]
and it is the algebraic reason that the Lie algebra $\mathfrak{so}(4)$ splits as $\su(2)\oplus\su(2)$. The splitting reappears in physics as the self-dual and anti-self-dual parts of a two-form, and as the left and right chiral halves of a Dirac spinor. For the octonions, which we meet in Interlude~A, left and right multiplications no longer commute, and the clean splitting above is replaced by triality.

\begin{litbox}
Everything in this chapter is classical: the relations \eqref{eq:hamilton} are Hamilton's \cite{Hamilton1844}, the isomorphism $S^3\cong\SU(2)$ and the double cover of $\SO(3)$ are standard \cite{ConwaySmith}, the fibration \eqref{eq:hopf} is due to Hopf \cite{Hopf1931}, and its many appearances in physics are surveyed by Urbantke \cite{Urbantke2003}. What the chapter adds is a point of view. We have stressed that $\Mat_2(\CC)$ carries two real forms with the determinant as a common quadratic form, and that the two-sphere of imaginary units is also the sphere of pure spin states. Both observations become load-bearing later, when the space of positive-energy states of a massive particle is identified with a hemisphere whose boundary is this sphere (Chapter~4). All theorems displayed above have been checked in the Lean proof assistant.
\end{litbox}

\section*{What this chapter passes on}
\begin{itemize}[nosep]
\item The norm identity $q\qbar q=|q|^2$ and its multiplicativity (\cref{thm:normmul}).
\item The faithful matrix model $\rho$ and $S^3\cong\SU(2)$ (\cref{thm:rho}).
\item Conjugation $x\mapsto qx\qbar q$ as rotations of $\RR^3$ (\cref{thm:conj}).
\item The two-sphere of imaginary units, pure spin states and (soon) light rays.
\item Left and right multiplications as commuting copies of $\su(2)$ in $\mathfrak{so}(4)$ (\cref{thm:commutant}).
\end{itemize}

\section*{Exercises}
\begin{exercise}
Using only \eqref{eq:hamilton} and associativity, show that $ij=k$, $jk=i$ and $ki=j$. Then verify that $\rho(i)\rho(j)=\rho(k)$ with the matrices of \eqref{eq:rho}, as \cref{thm:rho}(i) requires.
\end{exercise}
\begin{exercise}
Write $p=a+b\,i+c\,j+d\,k$ and $q=e+f\,i+g\,j+h\,k$ and expand $|pq|^2=|p|^2|q|^2$ to obtain Euler's four-square identity.
\end{exercise}
\begin{exercise}
Show from \cref{thm:rho}(iii) and (iv) that $\rho(q)^\dagger\rho(q)=|q|^2\mathbf 1$ for every quaternion, so that $\rho(q)/|q|$ is unitary whenever $q\neq0$.
\end{exercise}
\begin{exercise}
Let $q=\cos(\theta/2)+\sin(\theta/2)\,k$. Compute $q\,i\,\qbar q$ and $q\,j\,\qbar q$ and confirm that $x\mapsto qx\qbar q$ rotates the $(i,j)$-plane by the angle $\theta$. Check the result against \cref{thm:conj}.
\end{exercise}
\begin{exercise}
Verify directly from the matrices above that $R_i$ commutes with $L_i$ and with $L_j$, and that $R_i^{\mathsf T}=-R_i$.
\end{exercise}

\begin{checklist}
\chk{AlphaI:i_sq,AlphaI:j_sq,AlphaI:k_sq,AlphaI:ijk}{Hamilton's relations $i^2=j^2=k^2=ijk=-1$ (four statements).$^*$}
\chk{AlphaI:normSq_mul'}{Multiplicativity of the quaternion norm, $|pq|^2=|p|^2|q|^2$.$^*$}
\chk{AlphaI:rho_mul}{The matrix model preserves products: $\rho(pq)=\rho(p)\rho(q)$.$^*$}
\chk{AlphaCompleteB:rho_one}{$\rho(1)=\mathbf 1$.$^*$}
\chk{AlphaCompleteB:rho_injective}{$\rho(q)=0$ implies $q=0$, so $\rho$ is injective.$^*$}
\chk{AlphaCompleteB:rho_star}{$\rho(\qbar q)=\rho(q)^\dagger$.$^*$}
\chk{AlphaI:det_rho}{$\det\rho(q)=|q|^2$.$^*$}
\chk{AlphaCompleteB:rho_SU2}{A unit quaternion is sent to $\SU(2)$: $\rho(q)\rho(q)^\dagger=\mathbf 1$ and $\det\rho(q)=1$.$^*$}
\chk{AlphaI:hopf_re}{$\operatorname{Re}(qx\qbar q)=|q|^2\operatorname{Re}x$; conjugation preserves pure quaternions.$^*$}
\chk{AlphaI:hopf_norm}{$|qx\qbar q|^2=|q|^4|x|^2$.$^*$}
\chk{AlphaVII_Commutant:Rm_mem}{For pure $b$, $R_b$ is antisymmetric and commutes with $L_i$ and $L_j$.$^*$}
\chk{AlphaVII_Commutant:commutant_so4}{An antisymmetric $4\times4$ matrix commuting with $L_i$ and $L_j$ equals $R_b$ with $b$ read off from its first column.$^*$}
\end{checklist}
